% !TEX program = xelatex
% Related Literature + Research Motivation (LaTeX snippet)
% Style: "English journal style" narrative; citations assume natbib (e.g., \usepackage{natbib})

\section{Related Literature and Research Motivation}
\label{sec:lit_motivation}

\subsection{Economic shocks, opportunity costs, and crime/conflict}
\label{subsec:econ_shocks_crime}

A canonical starting point for understanding why adverse economic conditions may induce violence is the opportunity-cost framework in \citet{becker1968crime}, where lower legal earnings increase the relative payoff to illicit activities. Building on this logic, a central empirical challenge is that economic conditions and violence are jointly determined: conflict destroys production, while economic downturns may trigger conflict. \citet{miguel2004economic} offer a seminal solution by using rainfall variation as an instrument for income growth in Sub-Saharan Africa. Their estimates imply a large causal effect of negative income shocks on the onset of civil conflict, and their discussion emphasizes that downturns expand the pool of underemployed young men and lower the opportunity cost of joining armed groups. Although their setting concerns civil war rather than street crime, the paper establishes a widely used causal template---external shocks $\rightarrow$ economic distress $\rightarrow$ violence---that subsequent work adapts to more granular outcomes and spatial units.

The crime literature has long debated whether aggregate business cycles translate into local criminal activity, partly because both crime measurement and identification remain difficult (e.g., the ``Great Recession and crime'' debate summarized in \citealp{dixcarneirosoaresulyssea2018economic}). A key insight of recent research is to move away from contemporaneous correlations and instead exploit plausibly exogenous shocks to local labor demand. In this respect, \citet{dixcarneirosoaresulyssea2018economic} (and its earlier NBER version \citealp{dixcarneirosoaresulyssea2017economic}) provide one of the most influential designs linking trade-induced labor-market dislocations to violent crime. Exploiting Brazil's large unilateral tariff reductions in the early 1990s, they construct region-specific tariff shocks using initial industry composition and sector-level tariff changes. Their event-time estimates show that regions facing larger tariff cuts experience a pronounced but temporary increase in homicide rates after liberalization is completed, with limited evidence of differential pre-trends. Importantly, they go beyond reduced-form estimates and propose a framework that leverages distinct dynamic responses of candidate mechanisms (labor market conditions, public goods provision, and inequality) to bound the share of the crime effect attributable to labor-market deterioration. Their conclusions point strongly to labor-market channels, suggesting that employment and earnings declines account for the majority of the observed crime response. This paper is particularly close to our setting because (i) it operationalizes a trade shock as an arguably exogenous, place-varying exposure; (ii) it highlights the dynamic nature of adjustment (medium-run deterioration rather than instantaneous spikes); and (iii) it explicitly discusses mechanisms and policy-relevant margins.

While the above evidence is largely from Latin America and focuses on homicides, a rapidly growing set of studies documents that adverse trade shocks can generate broader ``social costs'' through crime, political polarization, and social unrest. For China---a context where conventional crime data are scarce---\citet{ma2025export} make a major contribution by constructing city-by-year crime measures from millions of court judgment documents and estimating the criminogenic effects of export slowdowns. Their identification relies on a Bartik/shift-share design: city exposure is a weighted average of foreign demand shifts by product, with weights given by the pre-period export composition. The paper documents robust increases in crime in cities more exposed to export contractions, with stronger effects in manufacturing-specializing areas and places with higher concentrations of young people and migrants. Mechanism evidence indicates that adverse export shocks reduce employment and migrant wages and worsen mental health, whereas alternative channels such as ``social stability'' spending play a smaller role. Relative to existing trade--crime work, this study is especially relevant to our project because it (i) provides China-specific evidence using an external-demand shock that is conceptually close to the U.S.--China trade war; (ii) uses text-based judicial data, offering a blueprint for measurement when police-record data are unavailable; and (iii) explicitly highlights migrant exposure and heterogeneity across local economic structures.

\subsection{Trade wars and tariff exposure: measuring shocks and interpreting incidence}
\label{subsec:tariff_exposure}

The 2018 tariff escalation between the United States and China has generated an active literature on tariff incidence, price pass-through, and distributional impacts. \citet{amiti2019impact} provide an accessible synthesis of early evidence, emphasizing that the burden of U.S. tariffs largely fell on U.S. importers and consumers through higher duty-inclusive prices rather than lower foreign exporter prices. \citet{fajgelbaum2020return} develop a comprehensive empirical and quantitative assessment of the ``return to protectionism'' in the United States, estimating large declines in targeted imports and exports and documenting near-complete pass-through of tariffs to import prices. Their general-equilibrium accounting highlights both aggregate real-income losses and spatial distributional consequences, since protected and retaliated-against sectors are geographically concentrated.

For empirical work that seeks to translate sector-level trade shocks into place-level treatment intensity, the dominant approach is the shift-share (Bartik-style) exposure design. In the classic ``China shock'' setting, \citet{autor2013china} treat commuting zones as local labor markets and measure import-exposure growth as the interaction of initial industry shares and national import growth, instrumenting U.S. imports from China using contemporaneous Chinese export growth to other high-income countries. Their results show sizable and persistent declines in manufacturing employment and earnings in exposed areas, as well as increased reliance on transfer programs. Although our empirical context and outcome differ, this literature clarifies why trade shocks can be viewed as plausibly exogenous conditional on initial specialization, and it motivates the construction of tariff-exposure measures that combine national/foreign shocks with pre-determined local shares.

At the same time, a growing econometric literature cautions that shift-share designs require careful inference and transparent identifying assumptions. \citet{goldsmithpinkham2020bartik} ``open the black box'' of the Bartik instrument and emphasize that identification typically comes from cross-sectional variation in exposure shares. This framing links shift-share designs to difference-in-differences logic: the key assumption is not that shares are uncorrelated with outcome levels, but that differential exposure to common shocks is orthogonal to differential trends absent the shocks. \citet{adao2019shiftshare} further demonstrate that conventional standard errors can severely over-reject because residuals may be correlated across regions with similar shares, even when spatial proximity is limited. They propose inference procedures that remain valid under arbitrary cross-regional residual correlation given independence of shifters across sectors. These methodological insights are directly relevant for our empirical strategy, which constructs a tariff-exposure index using a shift-share formula and estimates effects in a continuous-treatment DiD framework. In particular, our design inherits the strengths of exposure-based quasi-experiments---transparent shock construction and intuitive economic interpretation---but also requires robust inference and extensive placebo/pre-trend diagnostics to validate identification.

\subsection{Urban villages as ``arrival cities'': informal housing, migrant enclaves, and urban integration}
\label{subsec:arrival_city}

The above trade--crime literature largely treats local economies (cities, counties, commuting zones) as the relevant spatial unit. Yet, in rapidly urbanizing settings, economic shocks may be transmitted through heterogeneous micro-environments within the same city---especially through neighborhoods that concentrate vulnerable workers and informal employment. Our project draws conceptual guidance from \citet{saunders2010arrival}, who synthesizes ethnographic evidence from multiple continents to propose the ``arrival city'' as a key stepping-stone in the rural-to-urban migration pathway. In this framework, informal settlements and transitional neighborhoods perform a dual role. When ``healthy,'' they provide low-cost housing connected to the city, dense social networks that facilitate chain migration, and abundant opportunities for micro-entrepreneurship and informal employment. These features enable first-generation migrants to accumulate savings and convert them into upward-mobility investments (small business formation, housing assets, or children's education). When ``unhealthy,'' arrival neighborhoods suffer from weak connectivity, high effective living costs, limited space for informal entrepreneurship, and insufficient public services; in such environments, economic shocks can more easily translate into despair, social disorganization, and crime.

In China, urban villages (\emph{chengzhongcun}) are widely recognized as a quintessential form of ``arrival city'' that provides low-rent housing and dense migrant networks, but they are also highly contested in urban governance and redevelopment debates. \citet{wang2009urbanization} document how Shenzhen's rapid expansion left many original rural villages physically embedded in the city while retaining collective land institutions and informal housing construction. They emphasize that migrants, facing hukou-based exclusion from formal housing and high market rents, rely heavily on urban villages for affordable accommodation close to jobs. This line of work highlights that urban villages are not simply ``slums'' but are functional components of China's urbanization machine, providing both housing and employment opportunities for low-income migrants; it also cautions that large-scale redevelopment may neglect migrant welfare and disrupt informal labor markets.

Related research on housing and labor market segmentation underscores that migrants' residential sorting is systematically linked to labor-market outcomes. Using fine-grained census data for Shanghai, \citet{liu2016segmentation} show that rural migrants exhibit higher residential segregation than urban locals and tend to concentrate in low-skilled, low-paying jobs. They further provide suggestive evidence that migrant enclave residence can be associated with better employment outcomes for rural migrants, consistent with strong social networks and information channels within enclaves. Taken together, this literature implies that migrant neighborhoods may simultaneously (i) buffer shocks through networks and informal job matching, and (ii) amplify shocks if residents are disproportionately exposed to volatile, low-skill employment and limited social insurance.

Despite rich qualitative and descriptive accounts of urban villages, causal evidence linking external economic shocks to crime specifically within urban villages remains scarce. One reason is measurement: standard crime statistics are rarely available at neighborhood scale, and administrative boundaries are too coarse to align with the micro-spatial structure of urban villages. A second reason is conceptual: the ``arrival city'' framework predicts heterogeneous vulnerability depending on neighborhood ``health'' (connectivity, affordability, informal entrepreneurship space, and public service access), but these dimensions are difficult to quantify consistently across cities and over time. These gaps motivate new measurement strategies that combine text-derived crime events, geocoding, and remote-sensing-based mapping of the built environment.

\subsection{New data frontiers: judicial text, LLMs, and micro-spatial crime measurement}
\label{subsec:llm_crime_data}

Recent advances in ``text as data'' and machine learning have broadened the empirical frontier in settings with limited structured administrative data. In China, court judgment documents constitute a uniquely large corpus with standardized legal language and rich descriptions of incident facts, but they are not immediately usable for spatial analysis because key fields (time, location, offense category) are embedded in narrative text. \citet{zhang2025llm} demonstrate that large language models can extract structured spatiotemporal information from judgment documents at scale and, combined with geocoding APIs, can generate a national dataset with approximately million-level crime records at street or neighborhood granularity. This contribution is crucial not only for criminology and geography, but also for applied microeconomics, where credible identification increasingly relies on high-frequency, high-resolution outcome measures.

Complementary work in Chinese economics has begun to use judicial text to measure crime outcomes for policy evaluation. For instance, \citet{yi2023povertycrime} construct county-level criminal case measures from judgment documents and use a DiD design to assess whether China's targeted poverty alleviation reduces criminal crime growth, with mechanism evidence pointing to income growth and employment improvements. Although their setting and policy differ from ours, the paper illustrates both the feasibility and the pitfalls of using judicial text as a quasi-administrative outcome measure (coverage, timing, and classification issues), reinforcing the value of transparent construction and robustness checks.

Our project builds on these advances but targets a different empirical object: micro-spatial crime dynamics within urban villages in response to the 2018 U.S. tariff shock. We adapt LLM-based extraction and geocoding pipelines to our research context and integrate them with satellite-image-based deep-learning maps of urban villages to form a panel of grid cells within the Pearl River Delta. This integration matters because it aligns outcomes (crime incidents) with treatment exposure (tariff intensity via local industrial composition) and with theoretically motivated neighborhood attributes (built environment, connectivity, and public-service access) at a scale that conventional city-level studies cannot observe.

\subsection{Research gap and contribution}
\label{subsec:gap_contrib}

Synthesizing the above literature reveals a clear ``missing link'' between macro-level trade shocks and micro-level urban vulnerability. First, the trade--crime literature provides compelling evidence that adverse trade shocks can increase crime through labor-market distress, but it typically operates at coarse spatial units where within-city heterogeneity is averaged out. Second, the urban-village and ``arrival city'' literature provides a nuanced theory of migrant neighborhood functioning and vulnerability, but it is predominantly qualitative and rarely connected to sharp external economic shocks that enable causal inference. Third, even when researchers study crime in China, measurement constraints often preclude grid-level outcomes aligned with built-environment heterogeneity.

Our study aims to connect these strands. We treat the 2018 U.S. tariff escalation as a relatively ``clean'' external export shock to China's tradable sectors and measure differential exposure using a shift-share tariff-exposure index. We then estimate how crime evolves in urban villages relative to other urban spaces in a continuous-treatment DiD/event-study framework, emphasizing dynamic responses and pre-trend validation as in \citet{dixcarneirosoaresulyssea2018economic} and \citet{ma2025export}. Guided by the ``arrival city'' framework \citep{saunders2010arrival}, we further examine heterogeneity by neighborhood ``health''---proxied by primary education accessibility, housing affordability, and proxies for informal entrepreneurship opportunities (e.g., density and diversity of restaurants)—to test whether these attributes mitigate or amplify the criminogenic effects of trade shocks.

Conceptually, our contribution is to provide an explicit empirical bridge from a qualitative migration-and-neighborhood theory to a causal design in applied economics. Empirically, our contribution is to introduce a uniquely high-resolution dataset that combines (i) LLM-extracted, geocoded street crime events from judicial documents and (ii) remote-sensing-based urban-village maps, enabling causal identification at the micro-spatial scale of the built environment. Methodologically, our work brings recent best practices in shift-share and continuous-treatment inference (e.g., \citealp{goldsmithpinkham2020bartik,adao2019shiftshare}) to a setting where both the treatment and the outcome vary within cities, and where standard clustering choices can be nontrivial. Substantively, we speak to a broader debate on the social costs of globalization backlash and the urban governance of migrant neighborhoods: if urban villages function as ``shock absorbers'' of urbanization, understanding when they become ``powder kegs'' is essential for designing place-based mitigation policies and for building resilient ``arrival'' neighborhoods in the era of trade-policy uncertainty.

